

\section{\fabriq{Dejo esto para discutir: Diagonalization of the background mode equation}}
\label{app:diagonalization}

The linearized mode equation, Eq.~\eqref{eq:linearized_mode}, reads
%
\begin{equation}
    \partial_\tau \boldsymbol{\eta}_\kappa = \frac{1}{2\Omega}\mathcal{M}_\kappa(\tau)\,\boldsymbol{\eta}_\kappa,
    \label{eq:app_mode_eq}
\end{equation}
%
with $\mathcal{M}_\kappa(\tau) = A_\kappa + \cos\psi_1(\tau)\,B$ and
%
\begin{equation}
    A_\kappa = -\kappa^2\begin{pmatrix}\alpha & -1\\ 1 & \alpha\end{pmatrix}, \quad
    B = 2\begin{pmatrix}0 & 1\\ 0 & -\alpha\end{pmatrix}.
    \label{eq:app_A_B}
\end{equation}
%
$A_\kappa$ has eigenvalues $\Lambda_\pm = \kappa^2(-\alpha\pm i)$, eigenvectors $e_+=(1,i)^T$, $e_-=(1,-i)^T$. With $S=(e_+|e_-)$, $\boldsymbol{\eta}_\kappa = S\mathbf{w}$, Eq.~\eqref{eq:app_mode_eq} becomes
%
\begin{equation}
    \partial_\tau \mathbf{w} = \frac{1}{2\Omega}\big[\Lambda + \cos\psi_1(\tau)\,\tilde B\big]\mathbf{w},
    \label{eq:app_w_frame}
\end{equation}
%
with $\Lambda \equiv \mathrm{diag}(\Lambda_+,\Lambda_-)$ and $\tilde B \equiv S^{-1}BS$,
%
\begin{equation}
    \tilde B = \begin{pmatrix} i-\alpha & \alpha-i\\ \alpha+i & -(i+\alpha)\end{pmatrix}.
    \label{eq:app_Btilde}
\end{equation}

\subsection{Fourier content of the periodic background}
\label{app:fourier}

For a solution locked at rotation number $n$, $\psi_1(\tau)\approx n\tau+\psi_\text{c}+x\sin\tau$, $x\equiv\alpha h_\text{a}/\Omega$, so $\cos\psi_1(\tau)$ is $2\pi$-periodic and
%
\begin{equation}
    \cos\psi_1(\tau) = \sum_{m=-\infty}^{\infty} c_m\, e^{im\tau}.
    \label{eq:app_cos_fourier}
\end{equation}
%
Using the Jacobi--Anger identity as in Eq.~\eqref{eq:avg_walker},
%
\begin{equation}
    c_m = \tfrac12\Big[J_{m-n}(x)\,e^{i\psi_\text{c}} + J_{-m-n}(x)\,e^{-i\psi_\text{c}}\Big],
    \label{eq:app_cm}
\end{equation}
%
with $c_{-m}=c_m^{*}$. At $m=0$, using the locking condition $\epsilon \equiv h_0-n\Omega/\alpha=(-1)^nJ_n(x)\sin\psi_\text{c}$ and the stable branch $(-1)^nJ_n(x)\cos\psi_\text{c}>0$,
%
\begin{equation}
    c_0 = \sqrt{J_n(x)^2-\epsilon^2}\,,
    \label{eq:app_c0}
\end{equation}
%
real, non-negative, and vanishing at the step edges.

\subsection{Single-resonance (two-level) reduction}
\label{app:two_level}

In the interaction picture, $\hat{\mathbf{w}}\equiv e^{-\Lambda\tau/2\Omega}\mathbf{w}$, Eq.~\eqref{eq:app_w_frame} becomes
%
\begin{equation}
    \partial_\tau\hat{\mathbf{w}} = \frac{\cos\psi_1(\tau)}{2\Omega}
    \begin{pmatrix}
    \tilde B_{11} & \tilde B_{12}\,e^{-i\kappa^2\tau/\Omega}\\
    \tilde B_{21}\,e^{i\kappa^2\tau/\Omega} & \tilde B_{22}
    \end{pmatrix}\hat{\mathbf{w}}.
    \label{eq:app_interaction_picture}
\end{equation}
%
Keeping only the term stationary near $\kappa^2\approx m\Omega$ ($c_0$ on the diagonal, $c_m$ off-diagonal) and removing the residual phase via $\hat w_1=e^{-i\delta\tau/2}u_1$, $\hat w_2=e^{i\delta\tau/2}u_2$, $\delta\equiv\kappa^2/\Omega-m$, gives an autonomous $2\times2$ problem with eigenvalues
%
\begin{align}
    \lambda_\pm &= -\frac{\alpha c_0}{2\Omega} \nonumber\\
    &\quad\pm \frac{1}{2\Omega}\sqrt{(1+\alpha^2)|c_m|^2-(c_0+\delta\Omega)^2}.
    \label{eq:app_two_level_eigs}
\end{align}
%
This treats each resonance $m$ as isolated. Combining several by taking the union of their windows overestimates the true unstable region, since it discards coupling between near-degenerate channels -- removed next.

\subsection{General multi-harmonic secular expansion}
\label{app:general_hill}

Expand the solution of Eq.~\eqref{eq:app_w_frame} as a Floquet--Fourier series,
%
\begin{equation}
    w_1 = e^{\lambda\tau}\!\sum_{k} a_k e^{ik\tau}, \quad
    w_2 = e^{\lambda\tau}\!\sum_{k} b_k e^{ik\tau},
    \label{eq:app_floquet_ansatz}
\end{equation}
%
with $\lambda$ the Floquet exponent. Matching the coefficient of $e^{ik\tau}$ gives the exact recursion
%
\begin{align}
    (\lambda+ik)\,a_k &= \frac{\Lambda_+}{2\Omega}a_k \nonumber\\
    &\;+ \frac{1}{2\Omega}\sum_{m}c_m\big[\tilde B_{11}\,a_{k-m}+\tilde B_{12}\,b_{k-m}\big], \label{eq:app_recursion_a}\\
    (\lambda+ik)\,b_k &= \frac{\Lambda_-}{2\Omega}b_k \nonumber\\
    &\;+ \frac{1}{2\Omega}\sum_{m}c_m\big[\tilde B_{21}\,a_{k-m}+\tilde B_{22}\,b_{k-m}\big]. \label{eq:app_recursion_b}
\end{align}
%
This is the Hill's-method reformulation of Eq.~\eqref{eq:app_w_frame}, exact for unrestricted $k$, and equivalent to $\lambda\mathbf v=\mathcal H\mathbf v$ for $\mathbf v=(\dots,a_k,\dots,b_k,\dots)^T$.

\subsection{Truncated form and instability criterion}
\label{app:truncation}

Retaining sidebands $k=0,\dots,k_\text{max}$ (needing $c_m$ for $|m|\le2k_\text{max}$) gives a finite matrix $\mathcal H(\kappa)$ of size $2(k_\text{max}+1)$. Its blocks follow directly from Eqs.~\eqref{eq:app_recursion_a}--\eqref{eq:app_recursion_b}: with the $(k_\text{max}{+}1)\times(k_\text{max}{+}1)$ coupling matrix
%
\begin{equation}
    C_{kk'} \equiv c_{k-k'}, \qquad D_\pm \equiv \Lambda_\pm\mathbb{1} - 2i\Omega\,\mathrm{diag}(0,1,\dots,k_\text{max}),
    \label{eq:app_C_D_def}
\end{equation}
%
the Hill matrix is
%
\begin{equation}
    \mathcal H = \frac{1}{2\Omega}
    \begin{pmatrix}
    \tilde B_{11}C + D_+ & \tilde B_{12}C \\
    \tilde B_{21}C & \tilde B_{22}C + D_-
    \end{pmatrix}.
    \label{eq:app_Hill_matrix}
\end{equation}
%
Diagonalizing $\mathcal H(\kappa)$ gives the growth rate
%
\begin{equation}
    \lambda_\text{max}(\kappa) \equiv \max_i \operatorname{Re}\big[\mathrm{eig}\,\mathcal H(\kappa)\big],
    \label{eq:app_growth_rate}
\end{equation}
%
and the overall stability criterion at fixed $(h_0,h_\text{a})$ is
%
\begin{equation}
    \text{unstable} \iff \max_{\kappa>0}\lambda_\text{max}(\kappa) > 0.
    \label{eq:app_overall_criterion}
\end{equation}
%
This $\lambda$ relates to the monodromy Floquet multiplier via $\mu_\kappa=e^{2\pi\lambda}$, so $|\mu_\kappa|>1\iff\operatorname{Re}\lambda>0$, consistent with the criterion used in the main text.

\subsection{Three-sideband case ($k_\text{max}=2$)}
\label{app:kmax2}

For $k_\text{max}=2$, Eq.~\eqref{eq:app_C_D_def} needs only $c_0,c_{\pm1},c_{\pm2}$, giving the $3\times3$ block
%
\begin{equation}
    C = \begin{pmatrix} c_0 & c_{-1} & c_{-2}\\ c_1 & c_0 & c_{-1} \\ c_2 & c_1 & c_0 \end{pmatrix},
    \label{eq:app_C_block}
\end{equation}
%
with $D_\pm=\Lambda_\pm\mathbb 1-2i\Omega\,\mathrm{diag}(0,1,2)$ from Eq.~\eqref{eq:app_C_D_def}, assembled into the $6\times6$ matrix of Eq.~\eqref{eq:app_Hill_matrix}. This is the matrix diagonalized numerically in the main text; it reproduces both the fully-unstable $n=1$ tongue at weak drive and the reentrant stable window (``waist'') that opens once $h_\text{a}$ is large enough for the $m=0,1,2$ channels to hybridize.

\subsection{Consistency with the single-resonance limit}
\label{app:consistency}

Setting all entries of $C$ to zero except the pair coupling $k-k'=n$ reduces Eq.~\eqref{eq:app_Hill_matrix}, after the phase transformation of Sec.~\ref{app:two_level}, exactly to Eq.~\eqref{eq:app_two_level_eigs}. The multi-harmonic reduction thus recovers the leading-order result for well-separated resonances, departing from it precisely where neighboring resonances overlap.
